\documentclass[preview]{standalone}

\usepackage{amsmath}
\usepackage{amssymb}
\usepackage{stellar}
\usepackage{definitions}

\begin{document}

\id{lebesgue-measure-exercises}
\genpage

\section{Exercises}

\begin{snippetexercise}{lebesgue-measure-ex-1}{}
    Evaluate the Lebesgue measure of the following subset of \(\realnumbers\)
    \[
        \bigcup_{n=1}^\infty \left[
            n-\frac{1}{3^n}, n + \frac{1}{2^{n-1}}
        \right]
    \]
\end{snippetexercise}

\begin{snippetsolution}{lebesgue-measure-ex-1-sol}{}
    Let \[
        \Phi_n = \left[
            n-\frac{1}{3^n}, n + \frac{1}{2^{n-1}}
        \right]
    \]
    We have the measure of a single internal
    \[
        \mu(\Phi_n) = \left(n + \frac{1|}{2^{n-1}}\right)
        - \left(n-\frac{1}{3^n}\right) = \frac{1}{2^{n-1}} + \frac{1}{3^n}
    \]
    We will now study whether the intervals are disjoint.
    \begin{enumerate}
        \item \emph{Upper bound of \(\Phi_n\):} \[ U_n = n + \frac{1}{2^{n-1}} \]
        \item \emph{Lower bound of \(\Phi_n\):} \[ N_n = (n + 1) - \frac{1}{3^{n+1}} \]
    \end{enumerate}
    An overlap occurs if \(U_n > L_{n+1}\):
    \begin{align*}
        n + \frac{1}{2^{n-1}} &> (n + 1) - \frac{1}{3^{n+1}}  \\
        \frac{1}{2^{n-1}} + \frac{1}{3^{n+1}} &> 1
    \end{align*}
    We have the following first cases:
    \begin{itemize}
        \item For \(n=1\), we have \(\frac{1}{2^0} + \frac{1}{3^2} = \frac{10}{9} > 1\) (Overlap);
        \item For \(n=2\), we have \(\frac{1}{2^1} + \frac{1}{3^3} = \frac12 + \frac{1}{27} < 1\) (No overlap);
    \end{itemize}
    Since the terms in the inequality are strictly decreasing as \(n\) increases, there are no overlap for \(n \geq 2\).
    Thus, only \(\Phi_1\) and \(\Phi_2\) intersect. We have
    \begin{align*}
        \Phi_1 &= \left[1 - \frac13, 1 + \frac11\right] = \left[\frac23, 2\right] \\
        \Phi_2 &= \left[2 - \frac19, 2 + \frac12\right] = \left[\frac{17}{9}, \frac{5}{2}\right]
    \end{align*}
    and their intersection
    \[
        \Phi_1 \intersection \Phi_2 = \left[\frac{17}{9}, 2\right],
        \quad \mu(\Phi_1 \intersection \Phi_2) = 2 - \frac{17}{9} = \frac19
    \]
    Finally, the Lebesgue measure is given by
    \begin{align*}
        \mu\left(
            \bigcup_{n=1}^\infty \left[
                n-\frac{1}{3^n}, n + \frac{1}{2^{n-1}}
            \right]
        \right) &= \left(
            \sum_{n=1}^\infty \mu(\Phi_n)
        \right) - \mu(\Phi_1 \intersection \Phi_2) \\
        &= \left(
            \sum_{n=1}^\infty \frac{1}{2^{n-1}}
        \right) + \left(
            \sum_{n=1}^\infty \frac{1}{3^n}
        \right) - \frac19 \\
        &= \left(
            \sum_{n=0}^\infty \frac{1}{2^{n}}
        \right) + \left(
            \sum_{n=0}^\infty \frac{1}{3^n}
        \right) -\frac{1}{3^0} - \frac19 \\
        &= \frac{1}{1-\frac12} + \frac{1}{1-\frac13} - 1 - \frac19 = \frac{43}{18}
    \end{align*}
\end{snippetsolution}

\end{document}